\documentclass[a4paper,11pt]{article}

\newcommand{\TPnumero}{Project 2}
\newcommand{\macrosPath}{../../preamble/macros.tex}
% =========================================================
% Tutorial / lab preamble — Time Series Analysis module
% article class + fancyhdr + tcolorbox + listings (English)
% Author : Aymen Ben Brik (aymen.benbrik@esprit.tn)
% =========================================================

\usepackage[utf8]{inputenc}
\usepackage[T1]{fontenc}
\usepackage[english]{babel}
\usepackage{lmodern}
\usepackage{amsmath, amssymb, amsthm, mathtools, bm}
\usepackage{graphicx}
\usepackage{booktabs}
\usepackage{array}
\usepackage{multirow}
\usepackage{colortbl}
\usepackage{xcolor}
\usepackage{tikz}
\usetikzlibrary{positioning, arrows.meta, calc, shapes, fit, backgrounds, matrix}
\usepackage{listings}
\usepackage[most]{tcolorbox}
\usepackage{geometry}
\geometry{a4paper, margin=2cm, headheight=15pt}
\usepackage{fancyhdr}
\usepackage[hidelinks]{hyperref}
\usepackage{enumitem}

% --- Colours ---
\definecolor{espritBlue}{RGB}{30, 60, 110}
\definecolor{espritAccent}{RGB}{200, 80, 30}
\definecolor{boxEx}{RGB}{30, 60, 110}
\definecolor{boxQ}{RGB}{180, 120, 0}
\definecolor{boxC}{RGB}{30, 100, 60}
\definecolor{boxAtt}{RGB}{200, 80, 30}

% --- Listings ---
\definecolor{codebg}{RGB}{248,248,248}
\definecolor{codeframe}{RGB}{220,220,220}
\definecolor{keyword}{RGB}{0,82,155}
\definecolor{comment}{RGB}{88,110,117}
\definecolor{string}{RGB}{133,153,0}

\lstdefinestyle{pythonstyle}{
  language=Python,
  basicstyle=\ttfamily\small,
  keywordstyle=\color{keyword}\bfseries,
  commentstyle=\color{comment}\itshape,
  stringstyle=\color{string},
  backgroundcolor=\color{codebg},
  frame=single,
  rulecolor=\color{codeframe},
  frameround=tttt,
  breaklines=true,
  showstringspaces=false,
  tabsize=2
}
\lstdefinestyle{Rstyle}{
  language=R,
  basicstyle=\ttfamily\small,
  keywordstyle=\color{keyword}\bfseries,
  commentstyle=\color{comment}\itshape,
  stringstyle=\color{string},
  backgroundcolor=\color{codebg},
  frame=single,
  rulecolor=\color{codeframe},
  frameround=tttt,
  breaklines=true,
  showstringspaces=false,
  tabsize=2
}
\lstset{style=pythonstyle}

% --- Common macros (configurable path) ---
\providecommand{\macrosPath}{../../preamble/macros.tex}
\input{\macrosPath}

% --- Exercise counter ---
\newcounter{excounter}
\renewcommand{\theexcounter}{\arabic{excounter}}

\newtcolorbox[use counter=excounter]{exercise}[1][]{
  enhanced, breakable,
  colback=boxEx!5, colframe=boxEx, fonttitle=\bfseries,
  title={Exercise~\theexcounter\ifstrempty{#1}{}{ -- #1}},
  arc=2pt, boxrule=0.6pt
}
\newtcolorbox{question}[1][]{
  enhanced, breakable,
  colback=boxQ!5, colframe=boxQ, fonttitle=\bfseries,
  title={Question\ifstrempty{#1}{}{ -- #1}},
  arc=2pt, boxrule=0.5pt
}
\newtcolorbox{solution}[1][]{
  enhanced, breakable,
  colback=boxC!5, colframe=boxC, fonttitle=\bfseries,
  title={Solution\ifstrempty{#1}{}{ -- #1}},
  arc=2pt, boxrule=0.5pt
}
\newtcolorbox{warning}[1][]{
  enhanced, breakable,
  colback=boxAtt!5, colframe=boxAtt, fonttitle=\bfseries,
  title={Warning\ifstrempty{#1}{}{ -- #1}},
  arc=2pt, boxrule=0.5pt
}

% --- Headers / footers ---
\pagestyle{fancy}
\fancyhf{}
\fancyhead[L]{\textsc{\moduleNom} -- \auteurNom}
\fancyhead[R]{\TPnumero}
\fancyfoot[C]{\thepage}
\fancyfoot[L]{\auteurInstitution}
\fancyfoot[R]{\auteurEmail}
\renewcommand{\headrulewidth}{0.4pt}
\renewcommand{\footrulewidth}{0.2pt}

% --- Placeholder ---
\providecommand{\TPnumero}{Lab}


\title{Project 2: Decomposition of monthly souvenir sales\\
\large Trend, seasonality and multiplicative model selection}
\author{\auteurNom\\\auteurEmail\\\auteurInstitution}
\date{\anneeUniversitaire}

\begin{document}
\maketitle

\section*{Context}

The dataset \texttt{data/SouvenirSales.csv} contains the monthly
sales (in Australian dollars) of a souvenir shop on the beach of
Queensland (Australia), from January 1995 to December 2001. The series
exhibits both a strong upward trend and a marked seasonal pattern with
a Christmas spike — a textbook example of a \emph{multiplicative}
time series.

\section*{Deliverables}
\begin{itemize}
  \item A reproducible Python script (\texttt{project2\_<name>.py},
        seed $42$).
  \item A 4--6 page PDF report with figures, tables and discussion.
  \item Optional 5-minute oral defence.
\end{itemize}

% =========================================================
\begin{exercise}[Exploratory analysis]
\textbf{(1)} Load the data, parse the dates, plot $Y_t$ vs.\ time.

\textbf{(2)} Comment on:
\begin{itemize}
  \item the trend (is it roughly linear, exponential, or something else?);
  \item the seasonal amplitude (constant or growing with the level?);
  \item the presence of any obvious outliers.
\end{itemize}

\textbf{(3)} Compute and report descriptive statistics (mean, std, min,
max, skewness, kurtosis).
\end{exercise}

% =========================================================
\begin{exercise}[Additive vs.\ multiplicative]
\textbf{(1)} Apply an additive decomposition with
\texttt{statsmodels.tsa.seasonal\_decompose(y, period=12, model="additive")}
and plot the four panels.

\textbf{(2)} Same with \texttt{model="multiplicative"}.

\textbf{(3)} Compare the residuals of the two models. Which is more
homoscedastic?

\textbf{(4)} Conclude: which model is more appropriate for this series,
and why?
\end{exercise}

% =========================================================
\begin{exercise}[Log transform and additive workflow]
\textbf{(1)} Take $\log Y_t$. Plot the log-series.

\textbf{(2)} On $\log Y_t$, apply the full additive workflow:
\begin{itemize}
  \item estimate the trend (LOESS, frac = 0.3, or polynomial degree 2);
  \item detrend;
  \item estimate seasonality with seasonal averaging \emph{and} a
        cosine--sine model with one harmonic;
  \item compute the residual.
\end{itemize}

\textbf{(3)} Diagnose the residual: histogram, QQ-plot, time-vs-residual
plot. Is it stationary?
\end{exercise}

% =========================================================
\begin{exercise}[Forecasting]
\textbf{(1)} Split the data: train = 1995-01 to 2000-12 (72 months),
test = 2001-01 to 2001-12 (12 months).

\textbf{(2)} On the train set, fit a model
$\log Y_t = \beta_0 + \beta_1\,t + \sum_{j=1}^{2}[a_j \cos(j\omega t)
+ b_j \sin(j\omega t)] + \varepsilon_t$ ($\omega = 2\pi/12$).
Use OLS.

\textbf{(3)} Predict $\widehat{\log Y}_t$ for $t$ in the test range.
Convert back: $\widehat{Y}_t = \exp(\widehat{\log Y}_t)$.

\textbf{(4)} Compute the test MSE and MAPE. Plot observed vs.\ predicted
on the test period.
\end{exercise}

% =========================================================
\begin{exercise}[Bonus: residual modelling]
\textbf{(1)} Plot the autocorrelation function of the train residual
(\texttt{statsmodels.graphics.tsaplots.plot\_acf}).

\textbf{(2)} Are there significant autocorrelations beyond lag $0$?
This would suggest that the residual is not white noise — and
motivates the ARMA modelling of Chapter~5.
\end{exercise}

% =========================================================
\section*{Marking grid (out of 100)}

\begin{center}
\begin{tabular}{p{8cm}rl}
\toprule
\textbf{Criterion} & \textbf{Points} & \\
\midrule
Exploratory plots + descriptive statistics            & 15 \\
Additive vs.\ multiplicative comparison + discussion  & 20 \\
Log + full additive workflow + residual diagnostic    & 25 \\
Forecast on test set: train/test split, OLS, MAPE     & 20 \\
ACF of residual + ARMA motivation (bonus)             & 10 \\
Reproducibility + report quality                      & 10 \\
\midrule
\textbf{Total}                                        & \textbf{100} \\
\bottomrule
\end{tabular}
\end{center}

\bigskip
\hfill\auteurNom\par
\hfill\auteurEmail\par

\end{document}
